\scanpage{106}
\mspage{\deleted{\unsure{8}}55}
$f$ une $\overset{\deleted{\unsure{-}}}{W}$-fibration, et $p\in\overset{\deleted{\unsure{-}}}{W}$, alors $\underline{q\in\bar W}$, i.e. $\mathrm{hot}_W(q)$ iso. Pour ceci, nous interprétons la flèche $q$ dans $\mathrm{Hot}_W$ \deleted{\unsure{comme}} de la façon suivante.
Soit
\[
\unsure{\xi|\zeta}=\mathrm{hot}_{\unsure{Y|S},W}(X/Y),\qquad \xi'=p^*(\deleted{\unsure{Y}}\xi),
\]
je dis que $q$ s'identifie à la flèche canonique
\[
H_{\unsure{\bullet|\cdot}}^{\mathrm{HOT}_W}(Y',\xi')\longrightarrow H_{\bullet}^{\mathrm{HOT}_W}(Y,\xi)\quad\text{dans }\substack{\mathrm{Hot}_W\\ \Vert\\ \mathrm{HOT}_W(e)}\tag*{\deleted{\unsure{a}} (*)}
\]
où on note de plus que $\xi\in\mathrm{HOT}_W(\unsure{*|Y})$ est un coeff. \uline{loc. constant} sur $Y$.

Établissons d'abord cette interprétation de $q$. Notons que comme $X\in\underline{\mathrm{Fib}}_W Y$, l'image inverse $\xi'=p^*(\xi)$ n'est autre que $\mathrm{hot}_{Y',W}(X'=X\times_Y Y')$.
D'autre part, si \deleted{$g,g'$ sont} $g:Y\unsure{\relbar|\longrightarrow}e$, $g':Y'\longrightarrow e$ sont les flèches structurales, les deux \unsure{mbs|mbres} de (*) s'interprètent \inserted{\unsure{respt}} comme $g'_!(\xi')$ et $g_!(\xi)$ \deleted{Or le} (les $g'_!,g_!$ au sens du prédérivateur $\mathrm{HOT}_W$).
Or le calcul des $g'_!,g_!$ nous donne que les deux mbs s'identifient resp. à
\[
\mathrm{hot}_W(X'),\qquad\mathrm{hot}_W(X),
\]
et nous admettrons que la flèche \deleted{\unsure{du}}
\[
\mathrm{hot}_W(\unsure{X''|X'})\longrightarrow\mathrm{hot}_W(\unsure{X|x})
\]
n'est autre que $\mathrm{hot}_W(q)$. D'ailleurs, \deleted{\unsure{si}} \unsure{puisque} $X$ est $W$-fibré sur $S$, par définition pratiquement $\xi$ est dans $\mathrm{HOT}_W^{\mathrm{lc}}(S)$. Donc il nous suffira de prouver la
% archive sheet number: 106
